\chapter{SWOT ANALİZİ}
\thispagestyle{empty}

\begin{itemize}
    \item \textbf{Güçlü Yönler (Strengths):}
    \begin{itemize}
        \item Wikipedia'nın hiperlink yapısı, makaleler arası ilişkileri doğal bir graf olarak sunmaktadır; ayrıca bir graf inşa etmeye gerek yoktur.
        \item Kullanılan cümle gömme modelleri (bge-m3 ve gte-modernbert-base), büyük ölçekli metin derlemleri üzerinde ön eğitimli olduğundan ek alan adaptasyonu gerektirmeden yüksek kaliteli anlamsal temsiller üretebilmektedir.
        \item GAT'ın dikkat mekanizması, her bağlantının önemini otomatik olarak öğrenerek öneri kalitesini artırmaktadır.
        \item Wiki-CS \cite{mernyei2020wikics} gibi hazır ve temizlenmiş kıyaslama veri setlerinin mevcut olması, modelin hızlıca test edilmesine olanak tanımaktadır.
    \end{itemize}

    \item \textbf{Zayıf Yönler (Weaknesses):}
    \begin{itemize}
        \item Cümle gömme modellerinin girdi uzunluğu sınırı, uzun Wikipedia makalelerinin tamamını işlemeyi engellemektedir; metin kesme veya özetleme stratejileri gerekecektir.
        \item GNN ve gömme modellerinin birlikte çalıştırılması yüksek GPU belleği gerektirmektedir; bu durum donanım kısıtları nedeniyle darboğaz oluşturabilir.
        \item Proje gerçeklenmek için hem veri mühendisliği hem de model geliştirme gibi iki farklı uzmanlık alanına ihtiyaç duymaktadır.
    \end{itemize}

    \item \textbf{Fırsatlar (Opportunities):}
    \begin{itemize}
        \item WikiRecNet \cite{wikirecnet2020} çalışmasının Doc2Vec kullandığı metin temsili katmanının modern cümle gömme modelleri ile modernize edilmesi, doğrudan bir iyileştirme fırsatı sunmaktadır.
        \item LLM'lerin hiyerarşik sunum katmanı olarak kullanılması, literatürde henüz yaygın olarak ele alınmamış yeni bir araştırma alanıdır.
        \item Geliştirilen yöntem, bilgisayar bilimi dışındaki Wikipedia alanlarına ve farklı dillere genişletilebilir potansiyele sahiptir.
    \end{itemize}

    \item \textbf{Tehditler ve Riskler (Threats):}
    \begin{itemize}
        \item Wikipedia'nın sürekli değişen yapısı, veri setinin zamanla güncelliğini yitirebilmesi riskini taşımaktadır.
        \item Büyük ölçekli graf üzerinde GNN eğitimi sırasında bellek yetersizliği (OOM) sorunlarıyla karşılaşılabilir.
        \item Öneri kalitesinin nicel olarak değerlendirilmesi için uygun bir ground-truth veri setinin bulunmaması, değerlendirme sürecini zorlaştırabilir.
        \item LLM katmanının halüsinasyon (hallucination) riski, hiyerarşik sunumun güvenilirliğini olumsuz etkileyebilir.
    \end{itemize}
\end{itemize}